\begin{abstract}
\begin{quote}
\textbf{v3 full-pivot update (2026-04-16)}: the paper's main claim is now multi-tool selection via Q-side coverage. K-side results are retained only insofar as they explain why stationary spectral steering fails on this regime and why the ontology basis is still geometrically meaningful.
\end{quote}

Enterprise AI agents routinely select \textbf{multiple tools} per query: a news-music request should emit \texttt{NewsTool} and \texttt{MusicTool} in sequence. Existing K-side spectral steering methods — SEKA, AdaSEKA, and related stationary attention interventions — apply the same perturbation at every decoding step. Under autoregressive re-attention, this biases the model toward reusing the same facet rather than covering the remaining ones. We therefore study a different intervention axis: \textbf{Q-side coverage}, $\Delta_Q^{(t)} = -\beta \sum_{s<t} P_{f_s} q_t$, which subtracts from the step-$t$ query the components aligned with already-emitted facets in a per-head ontology basis $B_\mathrm{ont}$. The purpose of the operator is simple: after one tool has been emitted, suppress its facet so the next decoding step can attend to an unused facet.


On MetaTool Subtask4 (multi-tool, 2-tool ground truth, N=497), Q-coverage at $\beta = -0.1$ improves F1 by \textbf{+1.64pp} on Qwen2.5-7B-Instruct and \textbf{+0.40pp} on Llama-3.1-8B-Instruct over the no-steer baseline, with three-tier null-control direction specificity (+4.0pp real vs controls on Qwen). On the same benchmark, stationary K-side steering is consistently negative: our K-bias gives $-4.6$pp on Qwen and $-31.2$pp on Llama, while SEKA gives $-7.95$pp on the available Llama slice. The paper's central claim is therefore an \textbf{axis-separation claim}: stationary K-side steering is a single-selection tool, while Q-side coverage opens a multi-selection regime that the K-side family does not reach.


The same ontology basis $B_\mathrm{ont}$ remains useful on the K-side, but only as a \textbf{secondary result}. First, it defines a strong stability diagnostic: at $\alpha=0.3$, real $B_\mathrm{ont}$ preserves FC emission (F1 0.685) while random and feature-shuffled controls collapse to F1 0.000, a direction-specificity gap of \textbf{+68.5pp} on Subtask4 N=497. Second, it serves as a compression axis: OCQ (1-bit categorical on $B_\mathrm{ont}$ plus residual quantization) improves PPL by $-4.37$ over KIVI at 2-bit on Qwen2.5-7B WT2. These results matter because they show that the ontology basis is not an arbitrary feature basis even when the paper's primary novelty is on the Q side.


The main theoretical role of the ontology basis is therefore modest but important: it gives one shared geometric object through which the Q-side multi-tool result, the K-side stability result, and the compression side result can all be stated. Theorem 6.1 supplies the common attention-weighted perturbation language, and Theorem 6.19 packages the shared-basis observation. But the paper should be read first as a \textbf{multi-tool steering paper}, not as a broad unification paper with equal-weight contributions on three axes.


The core empirical message is compact. Q-coverage is positive on the benchmark where stationary K-side methods fail, and the ontology basis is specific rather than interchangeable. The secondary message is that the same basis has unusually strong K-side stability and a plausible compression use.
\end{abstract}
